\documentclass[11pt]{article}
\usepackage[margin=1in]{geometry}
\usepackage{amsmath, amssymb, amsfonts}
\usepackage{xcolor}
\usepackage{enumitem}
\usepackage{mathtools}
\usepackage{setspace}
\setstretch{1.15}

\definecolor{mygreen}{rgb}{0,0.55,0}

\newcommand{\defn}[1]{\textcolor{mygreen}{\textbf{Definition:} #1}}
\newcommand{\thm}[1]{\textcolor{mygreen}{\textbf{Theorem:} #1}}
\newcommand{\idea}[1]{\textcolor{mygreen}{\textbf{Key Idea:} #1}}

\begin{document}

\begin{center}
{\Large \textbf{STAT 545 Midterm Notesheet}}\\
\vspace{0.2cm}
Markov Chains --- Definitions, Techniques, and Worked Examples
\end{center}

\vspace{0.4cm}

%------------------------------------------------
\section*{1. Two-State Markov Chains}

\defn{A (time-homogeneous) Markov chain satisfies
\[
P(X_{n+1}=j \mid X_n=i, X_{n-1},\dots) = P_{ij}.
\]}

For a two-state chain $\{0,1\}$ with
\[
P =
\begin{pmatrix}
1-p & p\\
q & 1-q
\end{pmatrix},
\quad
\alpha = (\alpha_0,\alpha_1),
\]
we compute several useful quantities.

\subsection*{Two-Step Transitions}
\idea{Compute $P^2 = P \cdot P$ entrywise.}

\[
P^2 =
\begin{pmatrix}
(1-p)^2 + pq & p(2-p-q)\\
q(2-p-q) & (1-q)^2 + pq
\end{pmatrix}.
\]

\subsection*{Expectation of the State}
\idea{Since $X_n \in \{0,1\}$,
\[
E[X_n] = P(X_n = 1).
\]}

Thus,
\[
E[X_2] = (\alpha P^2)_1
= \alpha_0 P^2_{01} + \alpha_1 P^2_{11}.
\]

\subsection*{General $n$-Step Behavior}
\thm{For $p+q>0$, the $n$-step transition matrix is}
\[
P^n = \frac{1}{p+q}
\begin{pmatrix}
q + p c_n & p - p c_n\\
q - q c_n & p + q c_n
\end{pmatrix},
\quad c_n = (1-p-q)^n.
\]

Therefore,
\[
E[X_n]
= \frac{p}{p+q} + \left(\alpha_1 - \frac{p}{p+q}\right)c_n.
\]

\subsection*{Long-Run Limit}
\idea{If $|1-p-q|<1$, then $c_n \to 0$.}

\[
\lim_{n\to\infty} E[X_n] = \frac{p}{p+q}.
\]

\idea{This equals the stationary probability of state 1.}

\subsection*{Return Times}
\thm{For an irreducible Markov chain,
\[
E[T_i \mid X_0=i] = \frac{1}{\pi_i},
\]
where $\pi$ is the stationary distribution.}

Here,
\[
\pi = \left(\frac{q}{p+q}, \frac{p}{p+q}\right),
\quad
E[T_1 \mid X_0=1] = \frac{p+q}{p}.
\]

\subsection*{Closed-Form Expression for \texorpdfstring{$P^n$}{Pn}}

For the two-state Markov chain with transition matrix
\[
P =
\begin{pmatrix}
1-p & p\\
q & 1-q
\end{pmatrix},
\qquad p,q \ge 0,\; p+q>0,
\]
the $n$-step transition matrix admits the closed-form solution
\[
\boxed{
P^n
=
\frac{1}{p+q}
\begin{pmatrix}
q + p(1-p-q)^n & p - p(1-p-q)^n\\[6pt]
q - q(1-p-q)^n & p + q(1-p-q)^n
\end{pmatrix}
}
\]

\idea{Key structure:}
\begin{itemize}[leftmargin=1.2em]
  \item $(1-p-q)^n$ is the \emph{decay term} (second eigenvalue raised to power $n$).
  \item The constant terms correspond to the stationary distribution
  \[
  \pi=\left(\frac{q}{p+q},\frac{p}{p+q}\right).
  \]
  \item As $n\to\infty$, the exponential term vanishes whenever $|1-p-q|<1$.
\end{itemize}

\idea{Fast consistency check (exam trick):}
\begin{itemize}[leftmargin=1.2em]
  \item Rows sum to 1.
  \item Setting $n=0$ recovers $P^0=I$.
  \item Taking $n\to\infty$ gives rows equal to $\pi$.
\end{itemize}

%------------------------------------------------
\section*{2. Random Walks on Graphs}

\defn{A simple random walk moves uniformly to a neighboring vertex.}

\subsection*{Transition Matrices}
\idea{For vertex $i$,
\[
P_{ij} = \frac{1}{\deg(i)} \quad \text{if } i \sim j.
\]}

Always compute degrees first, then normalize rows.

\subsection*{Communication Classes}
\defn{States $i,j$ communicate if $i \to j$ and $j \to i$.}

\idea{Disconnected components $\Rightarrow$ multiple classes.}

\subsection*{Recurrence}
\thm{All states in a finite irreducible Markov chain are recurrent.}

\subsection*{Period}
\defn{The period of state $i$ is
\[
d(i) = \gcd\{n \ge 1 : P^n_{ii} > 0\}.
\]}

\idea{Bipartite graphs typically have period 2.}

\subsection*{Ergodicity}
\thm{A Markov chain is ergodic iff it is irreducible and aperiodic.}

\idea{Lazy random walks (self-loop probability $>0$) are always aperiodic.}

%------------------------------------------------
\section*{3. Absorbing Markov Chains}

\defn{A state is absorbing if $P_{ii}=1$.}

\idea{Reorder states so transient states come first.}

\[
P =
\begin{pmatrix}
Q & R\\
0 & S
\end{pmatrix}
\]

\subsection*{Fundamental Matrix}
\thm{The fundamental matrix is
\[
F = (I-Q)^{-1}.
\]}

\idea{$F_{ij}$ = expected number of visits to $j$ starting from $i$ before absorption.}

\subsection*{Expected Time to Absorption}
\[
E[T \mid X_0=i] = \sum_j F_{ij}.
\]

\subsection*{Absorption Probabilities}
\[
\text{Absorption probabilities} = FR.
\]

\subsection*{Limiting Distribution}
\thm{If all closed classes are aperiodic,
\[
\lim_{n\to\infty} P^n =
\begin{pmatrix}
0 & (I-Q)^{-1} R \Pi\\
0 & \Pi
\end{pmatrix},
\]
where $\Pi$ contains stationary distributions of closed classes.}

%------------------------------------------------
\section*{4. Estimating Transition Matrices}

\subsection*{Maximum Likelihood Estimation}
\defn{The MLE of $P$ is obtained by row-normalizing transition counts.}

\idea{Small sample sizes $\Rightarrow$ unreliable estimates.}

\subsection*{Laplace Smoothing}
\defn{Add 1 to each count:
\[
\hat p_i = \frac{N_i + 1}{\sum_j (N_j + 1)}.
\]}

\idea{Prevents zero probabilities.}

\subsection*{Dirichlet (Empirical Bayes) Smoothing}
\defn{Use a Dirichlet prior with mean $\mu$ and strength $\alpha$.}

Posterior mean:
\[
\tilde p_i = \frac{N_i + \alpha \mu_i}{N + \alpha}.
\]

\idea{Borrow strength from other rows (pooled information).}



%------------------------------------------------
\end{document}
